This benchmark is shaped by a few deliberate design choices. The first is the use of document-level \texttt{CELEX} alignment rather than finer paragraph-level alignment. That decision gives the benchmark a stable and auditable cross-lingual anchor, but it also makes relevance coarser than span-level supervision. The second is the separation between retrieval context and generation context, which reflects the practical observation that the best text for retrieval is not always the best text for producing focused supervision. The third is the use of same-language legal QA generation rather than English-first generation plus translation, since legal query quality depends strongly on language-specific phrasing and on selecting one decisive legal point.

The current construction also has clear limitations. QA-candidate filtering can behave unevenly across languages because formatting conventions and paragraph structure differ across the source material. Even after staged validation, some accepted questions may remain weaker than the ideal single-focus legal retrieval query, especially when they are still too clause-heavy or too close to the original legal phrasing. More broadly, document-level relevance cannot fully capture the finer-grained evidence that a legal user may care about inside a long act or annex.

At the same time, these limitations help define the benchmark's next development steps. Future work can refine the legal question-shape controls, study language-specific failure patterns more systematically, and extend the released paper with stronger empirical analysis of how different model classes behave on this dataset. This makes the discussion section less a post-hoc caveat list and more a roadmap for benchmark maturation.
